\documentclass{report}
\usepackage{amsmath,amssymb}
\setlength{\parindent}{0mm}
\setlength{\parskip}{1em}
\begin{document}
\begin{center}
\rule{6in}{1pt} \
{\large Umberto, U. Villa \\
{\bf Block AMG Preconditioners For Mixed Finite Element Discretization of Porous Media Flow Problems}}

Lawrence Livermore National Laboratory \\ P O Box 808 \\ L-560 \\ Livermore \\ CA 94551 \\ U S
\\
{\tt villa13@llnl.gov}\\
Ilya Lashuk\\
Panayot Vassilevski\end{center}

\newcommand{\curl}{\operatorname{curl}}
\renewcommand{\div}{\operatorname{div}}
\def\Nedelec{N\'ed\'elec}
\newcommand{\hypre}{{\em hypre }}

The Darcy and the Brinkman equations are two fundamental models
describing the dynamics of an inviscid (Darcy) or viscid (Brinkman) fluid
in a matrix of an inhomogeneous porous medium, alternating bubbles and
open channels. Typical applications of these models are in underground
water hydrology, petroleum industry, automotive industry, and biomedical
engineering. In these applications, the high variability in the PDE
coefficients, that may take extremely large or small values, negatively
affects the conditioning of the discrete problem which poses a
substantial challenge for developing accurate and efficient solvers.\\
In this talk, we consider mixed formulations of the Darcy and Brinkman
equations based on the de Rham sequence $H(\curl)$-$H(\div)$-$L^2$. We
first give a brief overview of our previous results for
their finite element discretizations with \Nedelec, Raviart-Thomas and
piecewise discontinuous elements and for the construction of respective
block-diagonal AMG preconditioners. The theoretical results are
illustrated with numerical experiments.
Then we outline some progress towards the construction of specialized
coarse space correction
exploiting an element-based algebraic multigrid approach (AMGe) aimed at
ensuring better robustness of the resulting preconditioners.\\
We also present an application of practical relevance in the field of
petroleum industry based on the popular SPE10 dataset, a challenging
benchmark for oil reservoir simulations.\\

This talk is based on a joint work with Ilya Lashuk and Panayot Vassilevski.


\end{document}
